\chapter{गुणसूत्र}\label{ch:g9:chromosomes}

\Cref{rem:g9:heredity-and-traits:question} ने निर्धारकों को \omterm{def:g6:cell-unit-of-life:parts}{केंद्रक} में
घेर लिया था। किसी \omterm{def:g5:first-look-at-cells:cell}{कोशिका} को ठीक पल पर पकड़ो, और \omterm{def:g6:cell-unit-of-life:parts}{केंद्रक} अपने पत्ते खोल
देता है: उसके भीतर की चीज़ें लिपटकर गिनने और तस्वीर खींचने लायक़ पिंडों
में बदल जाती हैं --- यानी \omterm{def:g9:chromosomes:chromosomes}{गुणसूत्रों} में। और उन्हें गिनना, शरीर की
\omterm{def:g5:first-look-at-cells:cell}{कोशिकाओं} में भी और \omterm{def:g8:sexual-reproduction:gametes}{युग्मकों} में भी, वह क़रीब-क़रीब सब कुछ समझा देता है जो
परिवार का एलबम माँग रहा था: आधे, जोड़े, फेंटाई, और वह पलटी जो तय करती है
कि शिशु लड़का होगा या लड़की।

\section{केंद्रक, रँगे हाथों पकड़ा हुआ}

\begin{definition}[गुणसूत्र]\label{def:g9:chromosomes:chromosomes}
\emph{गुणसूत्र}\index{गुणसूत्र} \omterm{def:g6:cell-unit-of-life:parts}{केंद्रक} में धागे जैसा एक पिंड है, जो
(अच्छी तरह रँगा हुआ, और बँटवारे के वक़्त) तब दिखता है जब \omterm{def:g6:cell-unit-of-life:parts}{केंद्रक} की
सामग्री ढुलाई के लिए लिपट जाती है। बँटवारों के बीच वही धागे ढीले और अनदेखे
पड़े रहते हैं --- यानी काम करने वाली हालत; और बँटवारे पर वे गाढ़े हो जाते
हैं --- यानी चलने वाली हालत। तस्वीर लेकर और छाँटकर किसी \omterm{def:g5:first-look-at-cells:cell}{कोशिका} के
गुणसूत्र उसका \emph{गुणसूत्र-चित्र}\index{गुणसूत्र-चित्र} बनाते हैं:
यानी पूरी गड्डी, बिछाकर रखी हुई।
\end{definition}

\begin{proposition}[इंसानी गिनतियाँ]\label{prop:g9:chromosomes:counts}
\omterm{def:g9:chromosomes:chromosomes}{गुणसूत्र-चित्र} तीन संख्याएँ देते हैं, और वे तीनों पूरा अध्याय ढोती हैं:
\begin{enumerate}
  \item इंसान के हर शरीर-कोशिका में \emph{46} \omterm{def:g9:chromosomes:chromosomes}{गुणसूत्र} होते हैं ---
        \omterm{def:g1:the-five-senses:sense}{त्वचा}, \omterm{prop:g7:digestion-and-nutrients:absorption}{यकृत}, \omterm{def:g8:nervous-communication:neuron}{न्यूरॉन}, सबमें बराबर
        (\cref{prop:g5:first-look-at-cells:growth} वाले बँटवारे उस
        गड्डी की हूबहू नक़ल हर \omterm{def:g5:first-look-at-cells:cell}{कोशिका} तक पहुँचाते हैं);
  \item बिछाकर रखने पर वे 46 छँटकर \emph{23 जोड़े} बनाते हैं, और हर
        जोड़े के दोनों साथी आकार तथा धारियों में मिलते-जुलते हैं;
  \item और \emph{\omterm{def:g8:sexual-reproduction:gametes}{युग्मक}} इस नियम को तोड़ते हैं: \omterm{def:g4:animal-reproduction:sexual}{अंडाणु} और \omterm{def:g8:reproductive-systems:male}{शुक्राणु}
        \emph{23} ढोते हैं --- यानी हर जोड़े से एक \omterm{def:g9:chromosomes:chromosomes}{गुणसूत्र}, और कोई
        जोड़ा पूरा नहीं।
\end{enumerate}
\end{proposition}
\admitted

\begin{example}[निषेचन का गणित]\label{ex:g9:chromosomes:arithmetic}
ये गिनतियाँ किसी उपपत्ति की तरह आपस में बैठ जाती हैं। \omterm{def:g8:sexual-reproduction:gametes}{युग्मक} 23-23 पर;
और \omterm{def:g5:flower-to-fruit:steps}{निषेचन} (\cref{prop:g8:sexual-reproduction:core}) उन्हें जोड़ देता
है: \omterm{prop:g8:sexual-reproduction:core}{युग्मनज} में $23 + 23 = 46$ --- यानी जोड़े बहाल, और हर जोड़े का एक
साथी माँ से, एक पिता से। इसके बाद संतान की हर शरीर-कोशिका दोनों जनकों का
योगदान ढोती है --- और ठीक यही \cref{rem:g4:animal-reproduction:mix} ने
एक पूरी अध्याय-पीढ़ी पहले बाहर से देखा था, बछेड़े और उसके चाँद के साथ।
\end{example}

\begin{remark}[युग्मकों को आधा क्यों करना ही है]\label{rem:g9:chromosomes:whyhalve}
मान लो वे न करते: 46 \omterm{def:g9:chromosomes:chromosomes}{गुणसूत्रों} वाले \omterm{def:g8:sexual-reproduction:gametes}{युग्मक} 92 \omterm{def:g9:chromosomes:chromosomes}{गुणसूत्रों} वाली संतान
बनाते, और फिर 184 \omterm{def:g9:chromosomes:chromosomes}{गुणसूत्रों} वाले नाती-पोते --- यानी गड्डी हर पीढ़ी पर
दोगुनी। \omterm{def:g8:sexual-reproduction:gametes}{युग्मकों} का वह ख़ास आधा करने वाला बँटवारा --- जिसमें हर एक को
\emph{हर जोड़े से एक साथी} मिलता है --- ही वह चीज़ है जो \omterm{def:g5:biodiversity:species}{प्रजाति} की
गिनती पीढ़ी-दर-पीढ़ी 46 पर टिकाए रखती है। और साथ में एक बोनस भी: किसी
\omterm{def:g8:sexual-reproduction:gametes}{युग्मक} को हर जोड़े का \emph{कौन-सा} साथी मिलेगा, यह जोड़े-दर-जोड़े अलग-अलग
तय होता है --- यानी एक ही जनक से $2^{23}$ मुमकिन गड्डियाँ, और
\cref{prop:g8:sexual-reproduction:shuffle} वाली फेंटाई आख़िरकार गिनी
हुई।
\end{remark}

\section{गुणसूत्र ही निर्धारक ढोते हैं}

\begin{proposition}[इसके सबूत कि यह गड्डी ही मायने रखती है]\label{prop:g9:chromosomes:evidence}
\omterm{def:g9:chromosomes:chromosomes}{गुणसूत्र} निर्धारकों की सवारियाँ \emph{हैं} --- यह मान्यता नहीं, सबूत
है:
\begin{itemize}
  \item वे ठीक वैसा ही बर्ताव करते हैं जैसा
        \cref{prop:g9:heredity-and-traits:conclusions} माँगती है:
        हर \omterm{def:g5:first-look-at-cells:cell}{कोशिका} में हाज़िर, \omterm{def:g8:sexual-reproduction:gametes}{युग्मकों} में आधे, \omterm{def:g5:flower-to-fruit:steps}{निषेचन} पर बहाल, और
        साथी-दर-साथी चुने जाने लायक़;
  \item \omterm{def:g9:chromosomes:chromosomes}{गुणसूत्र-चित्र} की दुर्घटनाएँ उस माल की गवाही देती हैं: जिस
        गड्डी में एक \omterm{def:g9:chromosomes:chromosomes}{गुणसूत्र} ज़्यादा या कम हो वह पूरे शरीर का विकास
        बदल देती है --- मिसाल के लिए, छोटे 21वें जोड़े की एक फ़ालतू
        प्रति उन \omterm{def:g9:heredity-and-traits:traits}{लक्षणों} का वह पहचाना हुआ जमघट बना देती है जिसे डाउन
        संलक्षण कहते हैं। एक धागे की गिनती चूकी, और बहुत-से \omterm{def:g9:heredity-and-traits:traits}{लक्षण} खिसक
        गए: यानी वे धागे ही \omterm{def:g9:heredity-and-traits:traits}{लक्षणों} के निर्देश ढोते हैं;
  \item और एक जोड़ा तो हर \omterm{def:g9:chromosomes:chromosomes}{गुणसूत्र-चित्र} में अपना माल ख़ुद बता देता है:
        यानी \emph{लिंग \omterm{def:g9:chromosomes:chromosomes}{गुणसूत्र}}।
\end{itemize}
\end{proposition}
\admitted

\begin{definition}[X, Y, और वह सिक्का]\label{def:g9:chromosomes:xy}
23वाँ जोड़ा दो रूपों में आता है: दो मिलते-जुलते \emph{X} \omterm{def:g9:chromosomes:chromosomes}{गुणसूत्र} ---
यानी किसी लड़की का \omterm{def:g9:chromosomes:chromosomes}{गुणसूत्र-चित्र}; या एक X और एक छोटा \emph{Y} --- यानी
कोई लड़का। और फिर वह आधा करना यह पलटी चलाता है: हर \omterm{def:g4:animal-reproduction:sexual}{अंडाणु} एक X ढोता है;
जबकि \omterm{def:g8:reproductive-systems:male}{शुक्राणु} X या Y ढोता है, आधा-आधा। X वाले \omterm{def:g8:reproductive-systems:male}{शुक्राणु} से \omterm{def:g5:flower-to-fruit:steps}{निषेचन} XX
बनाता है, और Y वाले से XY --- यानी दो में एक मौक़ा
(\cref{ch:g9:heredity-and-traits} वाली मौक़े की भाषा, आख़िरकार अपनी
कार्य-प्रणाली के साथ), जिसका फ़ैसला \omterm{def:g5:flower-to-fruit:steps}{निषेचन} के ऐन पल पर, इससे होता है कि
कौन-सा तैराक पहुँचा।
\end{definition}

\begin{omfigure}
\begin{tikzpicture}[scale=0.9]
  % mother
  \node[draw, thick, omDef, rounded corners=4pt, font=\small,
    align=center] (mo) at (1.2,4.2) {माँ\\ XX};
  \node[draw, thick, omDef, rounded corners=4pt, font=\small,
    align=center] (fa) at (7.8,4.2) {पिता\\ XY};
  % gametes
  \node[draw, omDef, circle, font=\footnotesize] (e1) at (1.2,2.4) {X};
  \node[draw, omDef, circle, font=\footnotesize] (s1) at (6.6,2.4) {X};
  \node[draw, omDef, circle, font=\footnotesize] (s2) at (9.0,2.4) {Y};
  \draw[->, thick] (mo) -- (e1);
  \draw[->, thick] (fa) -- (s1);
  \draw[->, thick] (fa) -- (s2);
  \node[font=\footnotesize, left] at (0.8,2.4) {हर अंडाणु};
  \node[font=\footnotesize, right, align=left] at (9.4,2.4)
    {शुक्राणु:\\ आधे X, आधे Y};
  % offspring
  \node[draw, thick, omThm, rounded corners=4pt, font=\small,
    align=center] (girl) at (3.4,0.4) {XX\\ एक लड़की};
  \node[draw, thick, omMeth, rounded corners=4pt, font=\small,
    align=center] (boy) at (5.8,0.4) {XY\\ एक लड़का};
  \draw[->, thick, omExo] (e1) -- (girl);
  \draw[->, thick, omExo] (s1) -- (girl);
  \draw[->, thick, omExo] (e1) -- (boy);
  \draw[->, thick, omExo] (s2) -- (boy);
  \node[font=\footnotesize, align=center] at (4.6,-0.5)
    {निषेचन पर, दो में एक मौक़ा};
\end{tikzpicture}

{\small 23वें जोड़े की वह सिक्का-पलटी: \omterm{def:g2:animals-grow-up:born}{अंडे} का X किसी X वाले या Y वाले
\omterm{def:g8:reproductive-systems:male}{शुक्राणु} से मिलता है --- और उसी मुलाक़ात के साथ \omterm{def:g9:chromosomes:chromosomes}{गुणसूत्र-चित्र} तय हो जाता
है।}
\end{omfigure}

\begin{example}[पुराने सवाल, छोटे जवाब]\label{ex:g9:chromosomes:answers}
यह गड्डी पुराने हिसाब चुकता कर देती है। कोई संतान लड़का होगी या लड़की,
इसका किसी भी जनक की इच्छाओं, खाने या मौसमों से कोई लेना-देना नहीं ---
वह तो वही पलटी है, जो \cref{def:g9:chromosomes:xy} की कार्य-प्रणाली में
चलती है (और पुरानी नाइंसाफ़ियों के लेखे के लिए: फ़ैसला करने वाला फ़र्क़
\emph{\omterm{def:g8:reproductive-systems:male}{शुक्राणु}} में सवार होता है)। बराबर आधे परिवार के एलबम का भी जवाब दे
देते हैं: माँ और पिता, दोनों एक-एक पूरी आधी गड्डी देते हैं --- 23 और 23
--- और \omterm{def:g9:heredity-and-traits:heredity}{आनुवंशिकता} का गणित दोनों वंशों को बराबर मान देता है।
\end{example}

\begin{method}[कोई गुणसूत्र-चित्र पढ़ना]\label{met:g9:chromosomes:read}
हाथ में \omterm{def:g9:chromosomes:chromosomes}{गुणसूत्र-चित्र} की कोई तस्वीर आए, तो:
\begin{enumerate}
  \item गिनो: किसी शरीर-कोशिका में 46 और किसी \omterm{def:g8:sexual-reproduction:gametes}{युग्मक} में 23 अपेक्षित
        हैं;
  \item जोड़े बनाओ: साथी आकार और धारियों से मिलाए हुए, यानी 22 आम
        जोड़े और साथ में 23वाँ;
  \item 23वाँ पढ़ो: XX या XY;
  \item और गिनतियाँ जाँचो: जिस जोड़े में एक या तीन सदस्य हों वह
        \cref{prop:g9:chromosomes:evidence} के अनुसार किसी खिसके हुए
        विकास को समझा देता है।
\end{enumerate}
\end{method}

\section{अभ्यास}

\begin{exercise}[$\star$]\label{exo:g9:chromosomes:1}
\omterm{def:g9:chromosomes:chromosomes}{गुणसूत्र} क्या है, और वह दिखता कब है?
\end{exercise}

\begin{exercise}[$\star$]\label{exo:g9:chromosomes:2}
इंसान की तीनों गिनतियाँ बताओ, और यह भी कि हर एक कहाँ लागू होती है।
\end{exercise}

\begin{exercise}[$\star$]\label{exo:g9:chromosomes:3}
\omterm{def:g5:flower-to-fruit:steps}{निषेचन} का गणित चलाओ, और बताओ कि वह क्या बहाल करता है।
\end{exercise}

\begin{exercise}[$\star$]\label{exo:g9:chromosomes:4}
\omterm{def:g8:sexual-reproduction:gametes}{युग्मकों} को गड्डी आधी क्यों करनी ही है? वरना जो भागती हुई बढ़त होती,
उसे दिखाओ।
\end{exercise}

\begin{exercise}[$\star$]\label{exo:g9:chromosomes:5}
23वें जोड़े के दोनों रूप क्या हैं, और उनमें से किसी एक का फ़ैसला किस जनक
का \omterm{def:g8:sexual-reproduction:gametes}{युग्मक} करता है?
\end{exercise}

\begin{exercise}[$\star$]\label{exo:g9:chromosomes:6}
21वें \omterm{def:g9:chromosomes:chromosomes}{गुणसूत्र} की एक फ़ालतू प्रति क्या पैदा करती है, और यह उन धागों के
बारे में क्या साबित करता है?
\end{exercise}

\begin{exercise}[$\star\star$]\label{exo:g9:chromosomes:7}
\cref{prop:g8:sexual-reproduction:shuffle} वाली फेंटाई इस अध्याय की
भाषा में ठीक-ठीक कहाँ होती है --- और एक जनक की फेंटाई की जगह कितनी बड़ी
है?
\end{exercise}

\begin{exercise}[$\star\star$]\label{exo:g9:chromosomes:8}
\cref{prop:g9:heredity-and-traits:conclusions} की हर माँग को \omterm{def:g9:chromosomes:chromosomes}{गुणसूत्रों}
के उस बर्ताव से मिलाओ जो उसे पूरा करता है।
\end{exercise}

\begin{exercise}[$\star\star$]\label{exo:g9:chromosomes:9}
``हर जन्म पर लड़का होने का मौक़ा दो में एक है, पहले चाहे जो हुआ हो।''
इसका बचाव कार्य-प्रणाली से करो --- लगातार तीन लड़कियों से कुछ भी क्यों
नहीं बदलता?
\end{exercise}

\begin{exercise}[$\star\star$]\label{exo:g9:chromosomes:10}
\cref{met:g9:chromosomes:read} इन पर चलाओ: (क) 46,XX वाली एक
शरीर-कोशिका; (ख) 23,Y वाला एक \omterm{def:g8:sexual-reproduction:gametes}{युग्मक}; (ग) 47 वाली एक शरीर-कोशिका,
जिसमें फ़ालतू 21वें जोड़े में है।
\end{exercise}

\begin{exercise}[$\star\star$]\label{exo:g9:chromosomes:11}
इतिहास ने राज्यों को बेटियाँ मिलने का दोष रानियों पर मढ़ा।
\cref{ex:g9:chromosomes:answers} से, युग्मक-दर-युग्मक, वह लेखा सुधारो।
\end{exercise}

\begin{exercise}[$\star\star\star$]\label{exo:g9:chromosomes:12}
शरीर की हर \omterm{def:g5:first-look-at-cells:cell}{कोशिका} पूरे 46 ढोती है --- फिर भी कोई \omterm{def:g8:nervous-communication:neuron}{न्यूरॉन} और कोई
त्वचा-कोशिका बिलकुल अलग होते हैं
(\cref{rem:g6:cell-unit-of-life:twoways})। इससे खड़ी होने वाली पहेली
बताओ, और उसके हल की दिशा भी (यानी अगले अध्याय का काम: किसी \omterm{def:g9:chromosomes:chromosomes}{गुणसूत्र} पर
\emph{पढ़ा} क्या जाता है)।
\end{exercise}

\section{समस्या: कोशिका-आनुवंशिकी की मेज़}

\begin{problem}[{सप्ताहांत समस्या --- एक ही रोशनी वाली मेज़ पर चार
गुणसूत्र-चित्र}]\label{pb:g9:chromosomes:1}
किसी अस्पताल की कोशिका-आनुवंशिकी प्रयोगशाला की सिखाने वाली मेज़:
\omterm{def:g9:chromosomes:chromosomes}{गुणसूत्र-चित्र} क --- 46 \omterm{def:g9:chromosomes:chromosomes}{गुणसूत्र}, 23वाँ जोड़ा XY; ख --- 46, 23वाँ जोड़ा
XX; ग --- 23 \omterm{def:g9:chromosomes:chromosomes}{गुणसूत्र}, और उनमें एक X; घ --- 47, 21वें जोड़े पर तीन
प्रतियों के साथ, और 23वाँ जोड़ा XX।

\medskip\noindent\textbf{भाग I --- पढ़तें।}
\begin{enumerate}
  \item क, ख और ग को \cref{met:g9:chromosomes:read} से पढ़ो: हर एक
        किसकी \omterm{def:g5:first-look-at-cells:cell}{कोशिकाएँ} हो सकती हैं?
  \item ग की गिनती कोई ग़लती नहीं है। वह किस तरह की \omterm{def:g5:first-look-at-cells:cell}{कोशिका} है, और
        उसका 23 बँटवारे का कौन-सा इतिहास दर्ज करता है?
  \item घ को पढ़ो: गिनती, जोड़ा, और वह विकास भी जिसका वह ऐलान करता है।
  \item क के 23वें जोड़े का फ़ैसला किस जनक के \omterm{def:g8:sexual-reproduction:gametes}{युग्मक} ने किया, और उस
        \omterm{def:g5:flower-to-fruit:steps}{निषेचन} पर संभावनाएँ क्या थीं?
\end{enumerate}

\medskip\noindent\textbf{भाग II --- गणित वाले सवाल।}
\begin{enumerate}[resume]
  \item एक प्रशिक्षु पूछता है कि किसी एक व्यक्ति की सारी शरीर-कोशिकाएँ
        \omterm{def:g9:chromosomes:chromosomes}{गुणसूत्र-चित्र} क से क्यों मिलती हैं। जवाब
        \cref{prop:g5:first-look-at-cells:growth} वाले नक़ल करने वाले
        बँटवारों से दो।
  \item कोई और पूछता है कि 46 \omterm{def:g5:first-look-at-cells:cell}{कोशिकाओं} वाले दो जनकों ने ग का 23 बनाया
        कैसे। उस ख़ास बँटवारे का नाम और उसका नियम बताओ।
  \item ग को किसी X वाले \omterm{def:g4:animal-reproduction:sexual}{अंडाणु} के साथ जोड़ो, और उसकी अपनी ही
        प्रयोगशाला के Y वाले पड़ोसी के साथ भी --- रुको: ग तो ख़ुद
        \emph{Y-रहित} है। बताओ कि ग का X किसी भी संतान के बारे में,
        जिसे बनाने में वह मदद करे, क्या पक्का कर देता है।
  \item एक जनक जितनी गड्डियाँ बाँट सकता है उन्हें निकालो:
        जोड़े-दर-जोड़े की आज़ादी बताओ और $2^{23}$ वाला नतीजा भी ---
        और फिर उसमें \omterm{def:g5:flower-to-fruit:steps}{निषेचन} का ऐसे दो बँटवारों को जोड़ना भी जोड़ो।
\end{enumerate}

\medskip\noindent\textbf{भाग III --- परामर्श का कमरा।}
\begin{enumerate}[resume]
  \item \omterm{def:g9:chromosomes:chromosomes}{गुणसूत्र-चित्र} घ वाली संतान के माता-पिता पूछते हैं, ``हमसे
        ग़लती क्या हुई?'' आधा करने की दुर्घटनाओं के बारे में
        कार्य-प्रणाली का ईमानदार जवाब दो --- और यह भी कि उस बच्चे के
        47 धागे उस कमरे के प्यार के बारे में क्या बदलते हैं और क्या
        नहीं। (जवाब सलाहकार की तरह दो, और जीव विज्ञान सीधा रखो।)
  \item तीन बेटियों वाले एक जोड़े को अगली बार बेटा होने की संभावना
        जाननी है। जवाब उसी पलटी से दो, और ``क़िस्मत के दौर'' वाली
        सहज-बुद्धि सुधारो।
  \item एक आगंतुक दावा करता है कि कोई परहेज़ शिशु का लिंग चुन सकता है।
        इस दावे की जाँच \cref{def:g9:chromosomes:xy} के सामने करो:
        फ़ैसला कहाँ होता है, और कोई परहेज़ पहुँच किस तक सकता है?
  \item मेज़ का समापन एक वाक्य में करो: वे चारों तस्वीरें मिलकर इस
        बारे में क्या साबित करती हैं कि \omterm{def:g9:heredity-and-traits:heredity}{आनुवंशिकता} सवारी किस पर करती
        है।
\end{enumerate}
\end{problem}
